\documentclass[11pt,a4paper]{article}
\usepackage[margin=22mm]{geometry}
\usepackage[T1]{fontenc}
\usepackage{lmodern}
\usepackage{microtype}
\usepackage{booktabs}
\usepackage{graphicx}
\usepackage{hyperref}
\usepackage{xcolor}
\hypersetup{colorlinks=true,linkcolor=blue,urlcolor=blue}
\graphicspath{{figures/}}

\title{Passive DS-TWR Speed-Limit Study}
\author{UWB Module ESP-IDF field validation}
\date{27 July 2026}

\begin{document}
\maketitle

\section{Purpose and method}

This experiment evaluates how far the custom receive-only Passive DS-TWR
protocol can be accelerated while retaining static accuracy. Four anchors
formed a surveyed \(3\,\mathrm{m}\times3\,\mathrm{m}\) square with
\(\pm2\,\mathrm{mm}\) placement tolerance. The tag was stationary at the
surveyed centre, \((1.5,1.5)\,\mathrm{m}\).

All measurements used firmware \texttt{137c2b5}. Positions were captured from
the lossless dashboard event stream rather than from periodic snapshots.
Raw AlgMin estimates and constant-velocity EKF estimates were retained
separately. The reported update rate counts complete position frames that
contain newly acquired UWB information.

\section{Results}

\begin{table}[ht]
\centering
\small
\begin{tabular}{lrrrrrr}
\toprule
Profile & pos/s & nominal & EKF RMSE & EKF P95 & raw RMSE & success\\
 & & pos/s & cm & cm & cm & \%\\
\midrule
Robust 16 ms & 49.3 & 62.5 & 1.73 & 2.90 & 2.05 & 98.2\\
Robust 13 ms & 50.6 & 76.9 & 1.89 & 3.19 & 2.21 & 93.8\\
Robust 10 ms & 61.5 & 100.0 & 1.95 & 3.23 & 2.25 & 92.1\\
Robust 10 ms, 0.9/0.9 ms & 62.8 & 100.0 & 2.58 & 4.14 & 2.82 & 91.1\\
Fast Star 10 ms & 56.7 & 100.0 & 3.83 & 6.15 & 4.12 & 92.9\\
\bottomrule
\end{tabular}
\caption{Controlled static measurements. Success is the aggregate
anchor-dialogue success rate reported by the firmware.}
\end{table}

\begin{figure}[ht]
\centering
\includegraphics[width=\linewidth]{measured_update_rate.png}
\caption{Measured position throughput. The dashed line marks the
85 positions/s target.}
\end{figure}

\begin{figure}[ht]
\centering
\includegraphics[width=\linewidth]{raw_vs_filtered_rmse.png}
\caption{The CV-EKF reduces dispersion at every tested operating point.}
\end{figure}

\begin{figure}[ht]
\centering
\includegraphics[width=0.84\linewidth]{speed_accuracy_tradeoff.png}
\caption{Speed--accuracy trade-off for the tested profiles.}
\end{figure}

\section{Interpretation}

The best fast operating point is \textbf{Robust Rotating with a configured
10 ms frame and 1.0/1.0 ms RESP/FINAL delays}. It delivered
\textbf{61.5 positions/s}, \textbf{1.95 cm 2-D RMSE}, and
\textbf{3.23 cm P95}. Robust 16 ms remains the most accurate profile at
1.73 cm RMSE, but the 10 ms profile adds 24.8\% update throughput for only
0.23 cm additional RMSE.

Fast Star did not improve speed in this geometry. Its asymmetric initiator
load produced 56.7 positions/s and 3.83 cm RMSE, so Robust Rotating is the
better optimization base.

The 0.9/0.9 ms delay experiment was also rejected. It improved throughput by
only 2.1\%, while RMSE rose from 1.95 to 2.58 cm and P95 from 3.23 to
4.14 cm. The live modules were therefore restored to the 1.0/1.0 ms control
timing.

\section{Why 85 positions/s was not reached}

The nominal 10 ms frame implies 300 directed slots/s. Anchor logs measured
approximately 196.7 slots/s, or 5.08 ms per effective slot. A complete
independent position frame uses three directed anchor dialogues, so
85 positions/s requires at least 255 successful slots/s before allowing for
rejections. Accepted position throughput must therefore improve by
approximately 38.2\%.

The present bottleneck is the radio transaction and scheduling path rather
than the dashboard or EKF. Re-solving after every observation could display
more than 85 updates/s, but consecutive results would reuse nearly all the
same measurements. Such a mode may reduce visual latency, but it must be
reported separately from independent frame throughput.

\section{Recommended next phase}

\begin{enumerate}
\item Instrument POLL TX, delayed RESP TX, FINAL TX/RX, CIA readout, and RX
re-arm with per-stage timestamps and failure counters.
\item Replace millisecond receive-loop scheduling with DW3000 delayed-TX and
delayed-RX deadlines in a non-blocking slot state machine.
\item Retain Robust Rotating at 1.0/1.0 ms as the control. Accept each timing
change only if RMSE and P95 remain inside the control confidence band.
\item Evaluate a rolling solve on every new observation as a distinct
low-latency display mode, exposing both solver updates/s and independent
frames/s.
\end{enumerate}

\section{Reproducibility}

The repository contains the raw JSONL event streams, status snapshots,
firmware timing logs, machine-readable metrics, and the analysis script
\texttt{tools/uwb\_passive\_speed\_analyze.py}.

\end{document}
